\chapter{Bitcoin Anchoring and the \texttt{MATH5} OP\_RETURN}
\label{ch:anchors}

\section{Motivation}
\label{sec:anc-mot}

A verification chain is only as trustworthy as the root that attests it.
Internal signatures prove the chain is self-consistent; they do not prove
\emph{when} or \emph{publicly} a given state was accepted. By anchoring
the Bifrost head (\cref{ch:bifrost}) to the Bitcoin ledger, MATHLIB5
gains a timestamp and a public, censorship-resistant root that no single
party controls.

\section{The Fingerprint}
\label{sec:anc-fp}

The root artifact is the SHA-256 of the published document (this PDF) and,
independently, the SHA-256 of the repository's verification receipt. Both
fingerprints are recorded in \texttt{fingerprint.txt}. A reader recomputes
the hash locally and compares; equality is necessary for authenticity.

\section{The Ed25519 Plasma Gate}
\label{sec:anc-ed}

Over the fingerprint, the Plasma Gate produces an Ed25519 signature with
the published public key (\texttt{pubkey.txt}). The signature witnesses
that the fingerprint was endorsed by the key holder at some time before
the anchor transaction. Verification (\texttt{verify\_anchors.py}) checks
\[
\text{Ed25519-Verify}\bigl(\text{pubkey}, \text{fingerprint},
\text{signature}\bigr) = 1.
\]

\section{The OP\_RETURN Payload}
\label{sec:anc-op}

The anchor is published as a Bitcoin \texttt{OP\_RETURN} output whose
payload is
\[
\texttt{MATH5:}\;\Vert\;\text{SHA-256}(\text{head}).
\]
The 6-byte \texttt{MATH5:} namespace distinguishes MATHLIB5 anchors from
other \texttt{OP\_RETURN} users; the 32-byte hash is the Bifrost head hash
(or, equivalently, the document fingerprint). The full payload is 38 bytes,
well within the 80-byte standard \texttt{OP\_RETURN} limit, and costs only
the dust transaction fee to embed.

\section{End-to-End Verification}
\label{sec:anc-e2e}

A verifier performs, in order:
\begin{enumerate}
  \item Fetch the Bitcoin transaction carrying the \texttt{MATH5:}
        \texttt{OP\_RETURN}.
  \item Confirm the prefix and recover the expected hash $h$.
  \item Hash the local document; require $\text{SHA-256}(doc) = h$.
  \item Require the Ed25519 signature over $h$ to validate under the
        published key.
\end{enumerate}
All four must hold for the document to be accepted as the anchored one.
The script \texttt{verify\_anchors.py} automates steps 2--4 and reports
\texttt{ALL CHECKS PASSED} on success.

\section{Lifetime and Rotation}
\label{sec:anc-life}

Each release produces a new anchor transaction; the chain of transactions
is itself a public timeline of the system's verification history. Key
rotation is supported by anchoring a successor public key; old anchors
remain valid for the intervals they covered, preserving the audit trail.

\section{Why Bitcoin}
\label{sec:anc-why}

Bitcoin is chosen not for its currency but for its adversarial,
proof-of-work timestamping: rewriting an anchor requires re-mining the
blocks above it, an economically bounded cost. Any ledger with comparable
finality would serve; Bitcoin is simply the most established such anchor
available without trusting a new party.
